% Manuscript fragment. Requires amsmath; no custom theorem environment.
% The earlier cusp-connection and uniform root-region theorems must
% precede this fragment. This file has not been LaTeX-compiled here.

\subsection{Monotonicity of the finite three-root width}

Let $c(\nu)=(t_*(\nu),\lambda_*(\nu),\mu_*(\nu))$ be the
validated cusp graph, $-29\leq\nu\leq0$. In the exact-cusp-centered
coordinates $s=t-t_*(\nu)$, $\ell=\lambda-\lambda_*(\nu)$ and
$m=\mu-\mu_*(\nu)$, write the two folds of the uniform local-region
theorem as $m_{\rm low}(\ell,\nu)<0<m_{\rm high}(\ell,\nu)$.
Set $W_\nu(\ell)=m_{\rm high}(\ell,\nu)-m_{\rm low}(\ell,\nu)$.

\paragraph{Finite-width monotonicity and nesting.}
For every $-29\leq\nu\leq0$ and $0<\ell\leq10^{-6}$,
\[
 \partial_\nu m_{\rm high}(\ell,\nu)<0,
 \qquad \partial_\nu m_{\rm low}(\ell,\nu)>0,
\]
and
\[
 0.0003\ell^{3/2}<-\partial_\nu W_\nu(\ell)
 <0.003\ell^{3/2}.
\]
Thus, as $\nu$ decreases, the three-root strips strictly expand by
inclusion after translating each section by its exact cusp. All root
counts refer to $|t-t_*(\nu)|\leq0.003$; no global real-axis claim is made.

\paragraph{Proof mechanism.}
Put $D_n=\partial_t^nF$ and, along the implicit fold parameterized by $s$,
\[
 B(s,\nu)=\frac{4\lambda_*'(\nu)D_2+D_6/4}{D_4}.
\]
Differentiation of $F=0$ at fixed $\ell$, using $F_t=0$ and the
hierarchy identities, gives $\partial_\nu m=B-\mu_*'$.
At the cusp $B(0)=\mu_*'$, $B_s(0)=B_{ss}(0)=0$ and
$B_{sss}(0)=-32(-D_3/D_4)'=-6\sqrt2\,C'(\nu)>0$.
Explicit fourth- and fifth-derivative bounds yield
$B_{sss}(s,\nu)>0$ uniformly for $|s|\leq3/4000$ on all 58
driver cells. Since $\ell\geq1.80s^2$ and
$1.80(3/4000)^2>10^{-6}$, this covers both folds for every stated
$\ell$. Three integrations give opposite strict signs of
$B(s,\nu)-B(0,\nu)$ on the two branches. The same bounds, together
with $1.80s^2\leq\ell\leq2.21s^2$, imply the displayed width-rate
inequalities. Complete remainder formulas, interval witnesses and the
independent rational recurrence are supplied in the verification appendix.

At the finite section $\ell=10^{-6}$, validated endpoint fold boxes give
\[
 2.49463101<100\left(\frac{W_{\nu_S}(10^{-6})}
 {W_0(10^{-6})}-1\right)<2.49463103.
\]
This finite-section percentage is distinct from the increase of the
limiting coefficient $C$, and is not asserted to be constant in $\ell$.
The S section fixes $\nu=\nu_S$ and allows $\lambda,\mu$ perturbations.
This differs from the original $\mu=0$ section used to locate S.
